\begin{abstract}
Floating-base robots must interact safely with humans while balancing under
rigid contact---a problem of \emph{interaction dynamics} at the arm
end-effector that is coupled, through the contact constraints, to the
whole-body balance task.
Classical whole-body control~(WBC) relies on fixed-gain impedance feedback
that accumulates steady-state error under sustained physical human--robot
interaction~(pHRI) forces.
This paper develops the \emph{interaction layer} of a floating-base WBC
hierarchy, extending the fixed-base two-layer Impedance
MPC~\cite{cao2026impedance} to this setting: a receding-horizon quadratic
program~(QP) that occupies the residual null space of the balance tasks and
predicts and rejects pHRI disturbances at the arm end-effector using a
Kalman-augmented disturbance state.
A contact-consistent feedback linearization---built on the contact-consistent
task-space inertia---reduces the end-effector plant to a double integrator
with a \emph{constant} state matrix within each contact mode, enabling offline
precomputation of the QP cost and ${\geq}1$\,kHz operation.
A covariance-inflation protocol preserves the disturbance estimate across
contact-mode switches, giving zero steady-state error under bounded constant
pHRI loads at fixed contact mode, and an Impedance Equivalence Theorem shows
the infinite-horizon limit recovers a classical task-space impedance whose
effective mass, damping, and stiffness adapt to posture and contact
configuration.
The surrounding balance hierarchy---a centroidal MPC over the single-rigid-body
dynamics and a priority-driven contact-consistent WBC---is the architecture
into which the interaction layer is designed to plug; in the simulations the
layer runs above a stabilizing balance controller.
In simulation on a 17-DOF biped and a Unitree G1 model, the primary benefit of
the contact-consistent formulation is \emph{stability} under dynamic contact
loading: a fixed-base impedance MPC of comparable gain diverges under periodic
contact shocks, whereas the contact-consistent layer stays stable and
sub-millimetre. In static stance the two are comparable, and the Kalman
augmentation provides a modest offset-free refinement.
\end{abstract}

\begin{IEEEkeywords}
Whole-body control, model predictive control, impedance control,
floating-base robots, physical human--robot interaction, operational space
formulation, Kalman filter, contact-consistent dynamics, legged
manipulation.
\end{IEEEkeywords}

% ----------------------------------------------------------
\section{Introduction}

Legged and floating-base robots must increasingly do more than locomote: they
must physically interact with people and their surroundings while keeping
balance. Safe interaction on such platforms is fundamentally a problem of
\emph{interaction dynamics}---the closed-loop relation between the arm
end-effector force and motion---coupled, through the contact constraints, to
the whole-body balance task.
These objectives are tightly coupled---arm motions shift the center of
mass~(CoM), changing contact force distribution, while ground reactions
propagate back through the body and appear as disturbances at the
end-effector.
Classical fixed-base impedance control~\cite{hogan1985impedance} and its
MPC extensions~\cite{cao2023passive,cao2026impedance} cannot address this
coupling because they assume the robot base is rigidly anchored.

The dominant paradigm for whole-body control of legged systems decouples
the problem into two layers: a centroidal MPC that optimizes ground
reaction forces~(GRFs) using a linearized single rigid-body dynamics~(SRBD)
model~\cite{dicarlo2018dynamic,kim2019highly}, and an inner WBC layer that
resolves these forces into joint torques via a prioritized
QP~\cite{bellicoso2016perception,koolen2016design}.
This architecture achieves remarkable locomotive agility---the MIT Cheetah
3~\cite{dicarlo2018dynamic} executes high-speed bounding and stair
climbing---but allocates 100\% of the robot's control authority to
locomotion and base-posture maintenance.
Any external arm interaction is treated as a disturbance to be suppressed,
not as a channel to be actively regulated.
A biped reaching to assist a human standing beside it, or a quadruped
manipulating a valve while maintaining stance, cannot be handled by these
frameworks with the compliance and zero-steady-state-error guarantees
required for safe pHRI.

Conversely, impedance MPC methods designed for fixed-base
manipulators~\cite{cao2026impedance,cao2023passive,haninger2022model,roveda2020model}
lack an unactuated base state, generalized-coordinate partitioning, or
contact-consistent mass inverses.
They assume an infinite-mass ground connection and cannot model the
propagation of foot contact forces to end-effector apparent inertia.
Deploying them directly on a floating-base platform produces steady-state
torque errors and potential instability during contact transitions.

The technical gap is the absence of a \emph{predictive},
\emph{disturbance-rejecting}, compliance-controlled manipulation layer that
natively integrates with an existing balance stack without degrading balance
performance. This paper closes that gap by adopting the fixed-base
two-layer Impedance MPC of~\cite{cao2026impedance} as the base architecture
and extending it to the floating-base, contact-switching setting, with the
following contributions:

\begin{enumerate}
\item \textbf{Contact-consistent residual plant} (Section~\ref{sec:osf}):
we show that the floating-base arm end-effector dynamics, after
priority-driven feedforward cancellation, reduce exactly to a linear double
integrator with a \emph{constant} discrete state matrix $A_d$ within each
contact mode. The sole modification from the fixed-base case is the
substitution of the contact-consistent mass inverse $\Mbar$ for the
free-space inverse $M^{-1}$ in the operational-space inertia $\Larm$.

\item \textbf{Contact-mode-indexed Impedance MPC}
(Section~\ref{sec:impedancempc}): a receding-horizon QP indexed to the
active contact mode selects the input matrix $B_d^{(m)}$ from a
precomputed library, maintaining the convex QP structure and enabling
${\geq}1$\,kHz update rates with an $N$-step prediction horizon.

\item \textbf{Kalman disturbance isolation}
(Section~\ref{sec:impedancempc}): the augmented Kalman state estimates external
pHRI forces, unmodeled leg-momentum variations, and SRBD approximation errors,
and propagates the estimate through the prediction horizon. At a fixed contact
mode this yields offset-free tracking (Theorem~\ref{thm:zss}); in simulation it
gives a modest steady-state refinement over the no-Kalman controller.

\item \textbf{Contact-transition protocol} (Section~\ref{sec:switching}):
a covariance-inflation rule at contact-mode switches preserves the
disturbance estimate across events, bounding the transient tracking error
as a function of Kalman re-convergence time.

\item \textbf{Impedance Equivalence Theorem} (Section~\ref{sec:theorem}):
the infinite-horizon limit of the proposed MPC recovers a classical
task-space impedance law with closed-loop effective mass
$M_{d,\text{eff}}=\Larm(q)$, damping $D_{\text{eff}}=\Larm(q)D_d$, and
stiffness $K_{\text{eff}}=\Larm(q)K_d$ all adapting to arm posture and
contact configuration through the contact-consistent inertia---without
requiring $\{M_d,D_d,K_d\}$ as online optimization variables. This
adaptation is a structural consequence of the contact-consistent feedforward
linearization and preserves the ${\geq}1$\,kHz update rate.

\item \textbf{Stability analysis} (Section~\ref{sec:stability}): we prove
asymptotic zero steady-state error for fixed contact mode and derive
explicit transient bounds across contact switches as a function of the
Kalman convergence time constant.
\end{enumerate}

The remainder of this paper is organized as follows.
Section~\ref{sec:related} surveys related work.
Section~\ref{sec:dynamics} develops the floating-base dynamics and contact
constraint formulation.
Section~\ref{sec:osf} derives the contact-consistent operational-space
plant.
Section~\ref{sec:impedancempc} presents the three-level Impedance MPC
architecture.
Section~\ref{sec:switching} describes contact-transition handling.
Section~\ref{sec:theorem} states and proves the Impedance Equivalence
Theorem.
Section~\ref{sec:stability} analyzes stability.
Section~\ref{sec:torque} gives the complete joint torque equation and an
architectural comparison with prior work.
Section~\ref{sec:simulation} reports simulation benchmarks on two robot
platforms.
Section~\ref{sec:conclusion} concludes.

% ----------------------------------------------------------
\section{Related Work}\label{sec:related}

\subsection{Centroidal MPC for Locomotion}

The dominant approach to predictive whole-body control projects the robot's
dynamics onto its centroidal frame.
Di~Carlo et~al.~\cite{dicarlo2018dynamic} introduced the convex SRBD
formulation for MIT Cheetah~3, linearizing the rotation dynamics about a
nominal pitch-roll and treating each support foot as a rigid contact.
The resulting time-varying linear system admits a QP solution at 40\,Hz for
the outer MPC, while an inner Whole-Body Impulse Control~(WBIC) layer maps
GRFs to joint torques at 500\,Hz.
This bi-level architecture is foundational to the present work; however,
Kim et~al.~\cite{kim2019highly} dedicate both layers entirely to
locomotion---external arm forces are filtered out through centroidal inertia
assumptions and rigid GRF assignment, and no mechanism exists for compliant
manipulation.

Bellicoso et~al.~\cite{bellicoso2016perception} proposed a hierarchical WBC
scheme for the ANYmal quadruped that solves a cascaded sequence of QPs to
track base and end-effector tasks simultaneously.
While effective for continuous trotting and terrain adaptation, the WBC
stack uses classical PD task objectives with instantaneous feedback
($N{=}1$); there is no receding-horizon mechanism to predict and pre-load
corrective torque before a disturbance develops.
Sustained contact forces produce the steady-state error quantified in
\eqref{eq:sse} below.

Koolen et~al.~\cite{koolen2016design} implemented a momentum-based
controller for Atlas that regulates centroidal momentum via a QP
distributing contact forces.
The formulation shares the centroidal perspective but does not include a
predictive loop for arm impedance and relies on high-gain stiffness to
suppress interaction errors.

Sleiman et~al.~\cite{sleiman2021unified} proposed a unified MPC framework
for whole-body dynamic locomotion and manipulation on legged robots,
simultaneously optimizing body posture, gait timing, and arm end-effector
trajectories over a receding horizon.
This represents the state of the art for integrated loco-manipulation and
is the closest prior work to the present architecture.
The key structural difference is that the arm task in~\cite{sleiman2021unified}
is formulated as a fixed-base QP subproblem: the task-space inertia uses
$M^{-1}$ rather than the contact-consistent $\Mbar$, and no integrating
disturbance state is included to guarantee zero steady-state pHRI error
under sustained contact.

\subsection{Trajectory Optimization and Nonlinear MPC}

Winkler et~al.~\cite{winkler2018gait} developed TOWR, a phase-based
trajectory optimizer that simultaneously synthesizes gait sequences,
foothold locations, and full-body motions over multi-second horizons using a
nonlinear program solved by interior-point methods.
The full rigid-body model captures leg-inertia coupling that SRBD ignores,
but restricts operation to ${\sim}20$--50\,Hz and offline planning.
Real-time pHRI management at 1\,kHz is outside the scope of this approach.

Grandia et~al.~\cite{grandia2023perceptive} extended TOWR with a real-time
model-predictive framework using sequential convex approximations, achieving
20\,Hz replanning for rough terrain.
While this improves responsiveness compared to pure offline planning, it
still lacks the integrating disturbance estimation required for zero
steady-state error under sustained contact.

\subsection{Impedance MPC for Manipulation}

Force control for robotic manipulators---encompassing impedance, admittance,
and hybrid position/force strategies---is surveyed comprehensively
in~\cite{villani2016force}.
The foundational impedance control law of Hogan~\cite{hogan1985impedance}
was unified with position and torque control into a passivity-preserving
framework by Albu-Sch\"{a}ffer, Ott, and
Hirzinger~\cite{albu2007unified}, providing the theoretical basis on which
predictive extensions build.

The present paper is a direct structural extension of a line of prior
work on saturated and predictive control.
Anti-windup designs for output tracking under actuator saturation and
constant disturbances~\cite{cao2004antiwindup}, and the associated
domain-of-attraction analysis~\cite{cao2002antiwindup}, established that an
integrating disturbance channel achieves zero steady-state tracking error
for fixed-base saturated linear systems---the foundational insight carried
forward here to the floating-base, contact-switching setting via the Kalman
augmented state.
The min--max MPC formulation for LPV systems~\cite{cao2005minmax} introduced
parameter-varying input matrices with input constraints, the direct
precursor to the contact-mode-indexed $B_d^{(m)}$ library of the present
work.
Building on these fixed-base results, Cao, Cheng, and
Li~\cite{cao2023passive} introduced a passive MPC framework for pHRI on
fixed-base manipulators in which the outer MPC optimizes impedance
parameters $\{M_d,D_d\}$ over a receding horizon; passivity is enforced via
a virtual energy tank.
Because impedance parameters enter nonlinearly into the prediction matrices,
iterative solvers are required and update rates are limited to 10--30\,Hz.
The two-layer Impedance MPC for pHRI~\cite{cao2026impedance} resolved this
rate limitation on fixed-base manipulators: an analytical feedforward
cancels gravity, Coriolis, and task-space inertia, reducing the residual
plant to a double integrator with a \emph{constant} state-transition
matrix, while the MPC decision variables become corrective \emph{forces}
$\Fmpc$ and a Kalman-augmented disturbance state provides a formal
zero-steady-state-error guarantee---sub-0.05\,mm under a sustained 15\,N
force on a 7-DOF manipulator, versus 44.8\,mm for classical impedance
control.
That framework is the direct base of the present paper, which extends it
from fixed-base manipulators to floating-base humanoids: the constant-$A_d$
structure is preserved by substituting the contact-consistent mass inverse
$\Mbar$ into the feedforward linearization, the Kalman state additionally
absorbs leg-momentum variations and SRBD approximation errors, and the
contact-mode-indexed structure of~\cite{cao2005minmax} handles stance-phase
transitions without sacrificing the precomputed cost matrix structure.

Haninger, Hegeler, and Peternel~\cite{haninger2022model} optimize force
references and impedance parameters jointly using stochastic MPC with
Gaussian Process models of task forces.
This provides complementary insights into uncertainty-aware impedance
shaping but does not address floating-base dynamics, underactuation, or
contact-consistent operational-space formulation.

\subsection{Operational Space Control and Floating-Base Inverse Dynamics}

Khatib~\cite{khatib1987unified} formulated the operational space control
framework, establishing task-space inertia, Coriolis compensation, and
dynamically-consistent pseudoinverses as the mathematical foundation for
task-level manipulation.
Sentis and Khatib~\cite{sentis2005synthesis} extended this to hierarchical
synthesis of whole-body behaviors, proving that priority-ordered null-space
projection guarantees non-interference between tasks---the SK05 law that
forms the backbone of Level~2 in the present architecture.
Righetti et~al.~\cite{righetti2011inverse} unified the floating-base
inverse dynamics perspective with external contact constraints, showing how
$\Mbar$ arises naturally from an orthogonal decomposition of the constrained
dynamics.
The present work builds on~\cite{khatib1987unified,sentis2005synthesis,righetti2011inverse}
by embedding a predictive MPC layer in the residual null space of the
floating-base contact-consistent hierarchy.

% ----------------------------------------------------------
\section{Floating-Base Robot Dynamics}\label{sec:dynamics}

\subsection{Generalized Coordinates}

A floating-base robot (humanoid, quadruped) has $n$ actuated joints and a
6-DOF unactuated base~\cite{righetti2011inverse}. The generalized
coordinates are:
\begin{equation}\label{eq:gencoord}
  q = \begin{bmatrix} q_b \\ q_j \end{bmatrix} \in \mathbb{R}^{n+6},
  \quad q_b \in SE(3),\; q_j \in \mathbb{R}^n
\end{equation}
where $q_b=(p_b,R_b)$ is the base position and orientation and $q_j$ are
the $n$ joint angles. The velocity vector is:
\begin{equation}\label{eq:vel}
  \dot{q} = \begin{bmatrix} v_b \\ \dot{q}_j \end{bmatrix} \in \mathbb{R}^{n+6},
  \quad v_b = \begin{bmatrix} \dot{p}_b \\ \omega_b \end{bmatrix}
\end{equation}

\subsection{Equations of Motion}

The floating-base equations of motion~\cite{righetti2011inverse} are:
\begin{equation}\label{eq:eom}
  M(q)\ddot{q} + C(q,\dot{q})\dot{q} + G(q) = S^\top\tau + J_c^\top(q)\lambda
\end{equation}
where $M(q)\in\mathbb{R}^{(n+6)\times(n+6)}$ is the positive-definite
inertia matrix; $h=C(q,\dot{q})\dot{q}+G(q)\in\mathbb{R}^{n+6}$ collects
Coriolis, centrifugal, and gravity terms; $S=[0_{n\times6},\,I_n]$ is the
selection matrix (the base has no direct actuation); $\tau\in\mathbb{R}^n$
are commanded joint torques; $J_c(q)\in\mathbb{R}^{n_c\times(n+6)}$ is the
contact Jacobian; and $\lambda\in\mathbb{R}^{n_c}$ are ground reaction
forces~(GRFs).

Partitioning \eqref{eq:eom} into base and joint blocks:
\begin{equation}\label{eq:partition}
  \begin{bmatrix}M_b & M_{bj} \\ M_{bj}^\top & M_j\end{bmatrix}
  \begin{bmatrix}\ddot{q}_b \\ \ddot{q}_j\end{bmatrix}
  +\begin{bmatrix}h_b \\ h_j\end{bmatrix}
  =\begin{bmatrix}0 \\ \tau\end{bmatrix}
  +\begin{bmatrix}J_{c,b}^\top \\ J_{c,j}^\top\end{bmatrix}\lambda
\end{equation}

\subsection{Rigid Contact Constraints}

A rigid contact at point $i$ enforces zero contact-point
velocity: $J_{c,i}(q)\dot{q}=0$~\cite{righetti2011inverse}.
Differentiating yields the acceleration-level constraint:
\begin{equation}\label{eq:contact_acc}
  J_{c,i}\ddot{q} = -\dot{J}_{c,i}\dot{q}
\end{equation}

Stacking all contacts and combining with \eqref{eq:eom}:
\begin{equation}\label{eq:KKT}
  \begin{bmatrix}M & -J_c^\top \\ J_c & 0\end{bmatrix}
  \begin{bmatrix}\ddot{q} \\ \lambda\end{bmatrix}
  =\begin{bmatrix}S^\top\tau-h \\ \gamma_c\end{bmatrix},
  \quad \gamma_c\triangleq-\dot{J}_c\dot{q}
\end{equation}

Solving for the GRFs:
\begin{equation}\label{eq:GRF}
  \lambda = \Lambda_c(q)\bigl(-\dot{J}_c\dot{q}-J_cM^{-1}(S^\top\tau-h)\bigr)
\end{equation}
where $\Lambda_c=(J_cM^{-1}J_c^\top)^{-1}$ is the contact-space inertia.

\subsection{Contact-Consistent Mass Inverse}

Define the contact-null-space projector~\cite{khatib1987unified,righetti2011inverse}:
\begin{equation}\label{eq:Pc}
  P_c = I - J_c^\top\Lambda_c J_c M^{-1}
\end{equation}
and the \textbf{contact-consistent mass
inverse}~\cite{khatib1987unified,righetti2011inverse}:
\begin{equation}\label{eq:Mbar}
  \Mbar = M^{-1}P_c^\top = M^{-1}(I-J_c^\top\Lambda_cJ_cM^{-1})
\end{equation}
$\Mbar$ replaces $M^{-1}$ in all operational-space formulas when contacts
are active. It projects out the contact-constraint subspace, ensuring task
forces do not violate kinematic contact constraints.

\subsection{Centroidal Dynamics}

The centroidal momentum $h_G=[k^\top,L^\top]^\top=A(q)\dot{q}\in\mathbb{R}^6$
aggregates the robot's linear and angular momentum about its
CoM~\cite{orin2013centroidal}. Differentiating:
\begin{equation}\label{eq:centroidal}
  \dot{h}_G = G_c(q)\lambda + \begin{bmatrix}mg \\ 0\end{bmatrix},
  \quad
  G_c = \begin{bmatrix}I_3 & I_3 & \cdots \\
                       (p_1{-}p_G)^\times & (p_2{-}p_G)^\times & \cdots\end{bmatrix}
\end{equation}

For the outer MPC, \eqref{eq:centroidal} is approximated by the
\textbf{SRBD} model~\cite{dicarlo2018dynamic}, treating the robot as a
lumped mass $m$ with constant inertia $I_G$. Linearizing about a nominal
orientation yields:
\begin{equation}\label{eq:SRBD}
  \dot{x}_c = A_c x_c + B_c(\{p_i\})u_c
\end{equation}
where $x_c\in\mathbb{R}^{12}$ is the centroidal state and $u_c$ collects
contact forces. The SRBD approximation error is
$O(m_{\text{leg}}/m_{\text{total}})^2$, acceptable for robots where leg
mass is below 20--30\% of total mass.

% ----------------------------------------------------------
\section{Operational Space Formulation and Whole-Body Hierarchy}
\label{sec:osf}

\subsection{Task-Space Dynamics}

For a task variable $x_i=\phi_i(q)\in\mathbb{R}^{m_i}$ with Jacobian
$J_i=\partial\phi_i/\partial q$, substituting \eqref{eq:eom} into the
task-space acceleration yields the contact-consistent task-space
dynamics~\cite{khatib1987unified,righetti2011inverse}:
\begin{equation}\label{eq:taskdyn}
  \Lambda_i(q)\ddot{x}_i + \mu_i = F_i
\end{equation}
where $\Lambda_i=(J_i\Mbar J_i^\top)^{-1}$ is the task-space inertia;
$\mu_i=\bar{J}_i^\top h - \Lambda_i\dot{J}_i\dot{q}$ collects projected
Coriolis, centrifugal, and gravity terms; $\bar{J}_i=\Mbar J_i^\top\Lambda_i$
is the dynamically-consistent pseudoinverse; and $F_i$ is the commanded
operational-space force.

\subsection{Hierarchical Task Synthesis (SK05)}

For $k$ tasks in priority order (Task~1 = highest), the Sentis--Khatib
law~\cite{sentis2005synthesis} synthesizes the control torque as:
\begin{equation}\label{eq:SK05}
  \tau = J_1^\top F_1 + \bar{N}_1^\top J_2^\top F_2
       + \bar{N}_{12}^\top J_3^\top F_3 + \cdots
       + \bar{N}_{1\cdots k}^\top\tau_{\text{null}}
\end{equation}
where $\bar{N}_{1\cdots i}=\prod_{j=1}^{i}(I-\bar{J}_jJ_j)$ is the
accumulated contact-consistent null-space projector.
The key property $J_j\bar{N}_{1\cdots i}=0$ for all $j\leq i$ guarantees
that the arm MPC layer (injected into the residual null space) cannot
destabilize the balance and contact-maintenance tasks above it.

Each task force $F_i$ in \eqref{eq:SK05} is typically a PD law:
\begin{equation}\label{eq:PD}
  F_i = \Lambda_i(\ddot{x}_{di}+K_{D,i}\dot{e}_i+K_{P,i}e_i)+\mu_i
\end{equation}
Under a persistent disturbance force $F_h$:
\begin{equation}\label{eq:sse}
  e_{\infty,i} = K_{P,i}^{-1}\Lambda_i^{-1}F_h \neq 0
\end{equation}
This residual steady-state error is the fundamental limitation of the SK05
PD law and the primary motivation for replacing the arm task slot with
Impedance MPC.

% ----------------------------------------------------------
\section{Three-Level Whole-Body Impedance MPC}\label{sec:impedancempc}

\subsection{Architecture Overview}

The control objectives are:
\begin{enumerate}
  \item maintain balance with contact forces inside friction cones;
  \item track an arm end-effector reference $p_d(t)$;
  \item reject pHRI forces with zero steady-state tracking error.
\end{enumerate}

These objectives are addressed by three nested control levels:

\textbf{Level~1 (Centroidal MPC, 40--100\,Hz):} plans CoM trajectory and
GRFs over a 500\,ms horizon via the SRBD model \eqref{eq:SRBD}.

\textbf{Level~2 (WBC Hierarchy, 500\,Hz):} resolves the Level~1 GRFs into
joint torques using the SK05 law \eqref{eq:SK05} for contact and balance
tasks, leaving the arm end-effector task slot open.

\textbf{Level~3 (Impedance MPC, $\geq$1\,kHz):} fills the arm slot with a
receding-horizon QP that predicts and rejects pHRI disturbances, replacing
the PD law \eqref{eq:PD} with a predictive force command.

\subsection{Contact-Consistent Residual Plant}

After Level~2 commits torques for contact maintenance and balance, the
residual arm end-effector dynamics are:
\begin{equation}\label{eq:plant}
  \Larm(q)\ddot{x}_{\text{arm}}+\mu_{\text{arm}} = F_{\text{arm}}+d_{\text{ext}}
\end{equation}
where $\Larm=(J_{\text{arm}}\Mbar J_{\text{arm}}^\top)^{-1}$ uses the
contact-consistent mass inverse \eqref{eq:Mbar},
$\mu_{\text{arm}}=\bar{J}_{\text{arm}}^\top h - \Larm\dot{J}_{\text{arm}}\dot{q}$,
and $d_{\text{ext}}$ is the pHRI wrench projected to arm task space.

\subsection{Contact-Consistent Feedforward and Horizon Freezing}

We execute an analytical operational-space feedforward:
\begin{subequations}\label{eq:feedforward}
\begin{align}
  F_{\text{arm}} &= \Larm(q)\ddot{p}_d + \mu_{\text{arm}} - \Fmpc \label{eq:ff_a}\\
  \tau_{\text{ff,arm}} &= S\,\Naug^\top J_{\text{arm}}^\top F_{\text{arm}} \label{eq:ff_b}
\end{align}
\end{subequations}

Level~2 updates $\Naug(q)$, $\Larm(q)$, and $\mu_{\text{arm}}$ at 500\,Hz.
Level~3 runs at $\geq$1\,kHz; on interleaved cycles the projector and
feedforward terms are held at their most recent Level~2 values.
The configuration changes by at most $\|\dot{q}\|\Delta t_2\approx0.002$\,rad
per Level~2 interval, so the frozen-matrix error is first-order small.

Substituting the feedforward \eqref{eq:ff_a} into the residual plant
\eqref{eq:plant} cancels the $\Larm\ddot p_d+\mu_{\text{arm}}$ terms, and with
the error convention $e_{\text{arm}}=x_{\text{arm}}-p_d$ (actual minus
desired) and $d(t)\triangleq\Larm^{-1}d_{\text{ext}}$ the residual tracking
error dynamics reduce to:
\begin{equation}\label{eq:errordyn}
  \ddot{e}_{\text{arm}} = -\Larm^{-1}(q)\,\Fmpc + d(t).
\end{equation}
The corrective force $\Fmpc$ thus enters with a restoring sign, so a positive
$\Fmpc$ drives $e_{\text{arm}}\to0$.

\begin{proposition}[Constant $A_d$ and Local LTI Horizon Freezing]
\label{prop:constant_Ad}
Within a fixed contact mode, by evaluating the configuration-dependent task
inertia at the current sampling instant $k$, the discrete-time state
transition matrix for $x_e(k)=[e_{\text{arm}}^\top,\dot{e}_{\text{arm}}^\top]^\top$
reduces to a constant linear system:
\begin{equation}\label{eq:Ad}
  A_d = \begin{bmatrix}I_3 & \Delta t\,I_3 \\ 0 & I_3\end{bmatrix},
  \quad
  B_d(k) = \begin{bmatrix}0 \\ -\Larm^{-1}(q_k)\,\Delta t\end{bmatrix}
\end{equation}
\end{proposition}

\begin{proof}
The continuous residual error dynamics \eqref{eq:errordyn} have state matrix
$A_c=\left[\begin{smallmatrix}0&I_3\\0&0\end{smallmatrix}\right]$, which is
nilpotent ($A_c^2=0$). The matrix exponential therefore terminates at first
order, $A_d=e^{A_c\Delta t}=I+A_c\Delta t$, which is \emph{exact} and
independent of configuration---the double-integrator structure in
\eqref{eq:Ad}. Evaluating $\Larm(q_k)$ at the current sample freezes
$B_d(k)=[0;\,-\Larm^{-1}(q_k)\Delta t]$ over the horizon, so $(A_d,B_d(k))$ is
LTI within the receding window; the configuration drift
$\|\dot q\|\Delta t_2$ between Level-2 updates enters only $B_d$ and is
first-order small.
\end{proof}

The constant $A_d$ property permits the QP cost matrix $H$ and its Cholesky
factor to be computed offline once per contact mode, reducing each online
MPC update to a matrix--vector multiply.

\subsection{Receding-Horizon QP}

The input matrix indexed to contact mode $m$ is:
\begin{equation}\label{eq:Bd}
  B_d^{(m)} = \begin{bmatrix}0 \\ -(\Lambda_{\text{arm}}^{(m)})^{-1}\Delta t\end{bmatrix}
\end{equation}

The $N$-step prediction matrix $\Gamma^{(m)}$ is constructed from
$(A_d,B_d^{(m)})$ by standard lifted-system expansion.
The receding-horizon QP is:
\begin{equation}\label{eq:QP}
  \min_{U}\;\tfrac{1}{2}U^\top H^{(m)}U + h^{(m)\top}U
  \quad\text{s.t.}\quad \|\Fmpc(k)\|_\infty\leq F_{\max}
\end{equation}
with $H^{(m)}=\Gamma^{(m)\top}\bar{Q}\Gamma^{(m)}+\bar{R}$ and
$h^{(m)}=\Gamma^{(m)\top}\bar{Q}x_{\text{free}}^{(m)}$.
The contact-mode index $m$ plays the role of the scheduling variable in the
LPV-MPC framework of~\cite{cao2005minmax}.
The box constraint $\|\Fmpc\|_\infty\leq F_{\max}$ is chosen so that arm
joint torques $\tau_{\text{arm}}=J_{\text{arm}}^\top\Fmpc$ remain below
hardware limits at all configurations within the operating
workspace~\cite{cao2002antiwindup}.
The QP~\eqref{eq:QP} is strictly convex and solved by OSQP~\cite{stellato2020osqp}
in $<0.1$\,ms for $N=20$.

\subsection{Kalman Disturbance Augmentation}

The disturbance $d(t)$ in \eqref{eq:errordyn} captures: (i)~external pHRI
forces; (ii)~unmodeled contact reactions from the feet; and (iii)~SRBD
approximation error.
Augmenting with an integrating disturbance state $\dhat\in\mathbb{R}^3$:
\begin{equation}\label{eq:augstate}
  \begin{bmatrix}x_e(k{+}1)\\\dhat(k{+}1)\end{bmatrix}
  =\underbrace{\begin{bmatrix}A_d & B_d^{(m)} \\ 0 & I\end{bmatrix}}_{A_{\text{aug}}}
  \begin{bmatrix}x_e(k)\\\dhat(k)\end{bmatrix}
  +\begin{bmatrix}B_d^{(m)}\\0\end{bmatrix}\Fmpc(k)+w(k)
\end{equation}
with $w\sim\mathcal{N}(0,Q_w)$ and $v\sim\mathcal{N}(0,R_v)$.
Unlike the strict random-walk model of~\cite{cao2026impedance}, in which
process noise drives only the disturbance state, $Q_w$ here retains a
small block on $x_e$---two orders of magnitude below the $\dhat$
block (Section~\ref{sec:simulation})---to regularize the filter against
the frozen-matrix and SRBD residuals of the floating-base setting; the
offset-free property is unaffected, as it follows from the integrating
structure of $A_{\text{aug}}$, not from the noise partition.
The Kalman gain $K_f$ is computed offline from the steady-state discrete
algebraic Riccati equation.
The integrating structure of $A_{\text{aug}}$ guarantees $\dhat(k)\to d$
for any bounded constant disturbance~\cite{cao2004antiwindup}.

The free-response prediction fed into the QP is:
\begin{equation}\label{eq:freeresponse}
  x_{\text{free}}^{(m)} = A_d^N x_e + \sum_{j=0}^{N-1}A_d^j B_d^{(m)}\dhat
\end{equation}
This causes the optimizer to pre-load corrective force before the
disturbance fully manifests at the end-effector.

% ----------------------------------------------------------
\section{Contact Transition Handling}\label{sec:switching}

\subsection{Contact-Mode Switch Protocol}

When a foot lifts off or touches down, $J_c$ changes discontinuously;
consequently $\Mbar$, $\Larm$, and $B_d$ all jump.
At contact switch time $t_s$, the following protocol is executed in order:
\begin{enumerate}
  \item Recompute $\Mbar$ with the new $J_c$.
  \item Recompute $\Larm$; select new mode index $m_{\text{new}}$; load $B_d^{(m_{\text{new}})}$.
  \item \textbf{Covariance inflation:} $P_{\text{aug}}\leftarrow\alpha P_{\text{aug}}$,
        $\alpha\in[3,5]$, to allow rapid re-estimation of the disturbance.
  \item \textbf{Hold $\dhat$:} do not reset to zero---balance-related
        disturbances persist across contact transitions.
\end{enumerate}
The Kalman filter re-converges within approximately $5\tau_{\text{Kalman}}$
samples, after which the disturbance estimate is again accurate.

\subsection{Contact-Mode-Indexed Matrix Library}

For a robot cycling through $K$ repeating contact modes, precompute
$B_d^{(m)}$ and $\Gamma^{(m)}$ for each mode.
Since $A_d$ is constant (Proposition~\ref{prop:constant_Ad}), $A_d^N$ is
precomputed once globally; mode-switch overhead is a pointer swap and
$\Gamma^{(m)}$ reload---under 0.1\,ms.

% ----------------------------------------------------------
\section{Impedance Equivalence Theorem}\label{sec:theorem}

\begin{theorem}[Whole-Body Impedance Equivalence]\label{thm:equivalence}
Under rigid contacts, fixed contact mode, and no disturbance, the
unconstrained infinite-horizon Impedance MPC (Level~3, $N\to\infty$)
renders the classical task-space impedance law, for the rendered impedance
gains $(K_d,D_d)$:
\begin{equation}\label{eq:impedance}
  \Larm(q)\ddot{e}+\Larm D_d\dot{e}+\Larm K_d e = F_h
\end{equation}
with effective configuration-adaptive mass $M_{d,\emph{eff}}=\Larm(q)$
depending on both arm posture and contact configuration.
\end{theorem}

\begin{proof}
In the infinite-horizon unconstrained limit with $\dhat=0$, the QP reduces to
the static discrete-LQR (Riccati) feedback $\Fmpc=K_\infty x_e$. Factoring the
Riccati gain through the task inertia, write
$\Fmpc=\Larm(q)\,(K_d\,e+D_d\,\dot e)$; this defines the acceleration-level
gains $(K_d,D_d)$. Substituting into \eqref{eq:errordyn} with
$d=\Larm^{-1}F_h$ gives
$\ddot e=-(K_d e+D_d\dot e)+\Larm^{-1}F_h$, and premultiplying by $\Larm$
yields \eqref{eq:impedance}. Hence the effective mass, damping, and stiffness
are $M_{d,\text{eff}}=\Larm(q)$, $D_{\text{eff}}=\Larm D_d$, and
$K_{\text{eff}}=\Larm K_d$.
The gains $(K_d,D_d)$ are the Riccati image of the QP weights $(Q,R)$, \emph{not}
the weights themselves, and the cheap-control limit $R\to0$ does not equal
$(K_d,D_d)$ (cf.\ the base framework~\cite{cao2026impedance}).
The contact-consistent $\Mbar$ in $\Larm$ makes the effective mass adapt to
both the arm joint configuration and the active contact footprint.
\end{proof}

Theorem~\ref{thm:equivalence} establishes that Impedance MPC generalizes
classical impedance control: the finite-horizon constrained QP adds
predictive disturbance rejection, constraint enforcement, and contact-mode
adaptation while reducing to the classical law in the limit.
It also gives the impedance gains $K_d$, $D_d$ a physical reading as the
\emph{target} stiffness and damping.
In the simulations of Section~\ref{sec:simulation},
the rendered impedance is $(K_d,D_d)=(6\times10^4\,I_3,\;60\,I_3)$; at the
nominal arm configuration ($\Larm\approx0.20\,I_3$\,kg) this yields
$K_{\text{eff}}=\Larm K_d\approx1.2\times10^4$\,N/m and
$D_{\text{eff}}=\Larm D_d\approx12$\,Ns/m, consistent with compliant pHRI
operation.

Notably, all three closed-loop impedance parameters---$M_{d,\text{eff}}$,
$D_{\text{eff}}$, and $K_{\text{eff}}$---adapt automatically as the arm
configuration and contact state change, while the cost weights $K_d$ and
$D_d$ remain fixed design parameters.
This adaptation is a structural by-product of $\Larm(q)$ and incurs no
additional solver cost, preserving the ${\geq}1$\,kHz control rate.

% ----------------------------------------------------------
\section{Stability Analysis}\label{sec:stability}

\subsection{Zero Steady-State Error Under Fixed Contact Mode}

\begin{theorem}[Zero Steady-State Error]\label{thm:zss}
Suppose the disturbance $d(t)$ in \eqref{eq:errordyn} satisfies
$\|d\|_\infty\leq\bar{d}<\infty$ and is asymptotically constant.
Under fixed contact mode, the Kalman-augmented closed-loop system with the
Level~3 QP feedback is asymptotically stable, and
$\lim_{k\to\infty}\|e_{\emph{arm}}(k)\|=0$.
\end{theorem}

\begin{proof}
The claim combines regulator stability with offset-free disturbance rejection;
these rest on two \emph{distinct} structural properties, which the augmented
model separates.

\emph{Regulator (control) side.} The \emph{un-augmented} pair
$(A_d,B_d^{(m)})$ is controllable: $B_d^{(m)}=[0;\,-\Larm^{-1}\Delta t]$ has
full column rank whenever $\Larm^{-1}$ is nonsingular, and then
$[\,B_d^{(m)}\;\; A_dB_d^{(m)}\,]$ has rank $6$. Two singularities can violate
this---(i)~\emph{kinematic arm singularities} (e.g., a fully-extended elbow)
where $J_{\text{arm}}$ drops rank, and (ii)~\emph{contact-state singularities}
where the platform loses all ground contact---so the result is stated on a
compact, singularity-free subset $\mathcal{W}$. On $\mathcal{W}$ the
infinite-horizon LQR gain $K_\infty$ places the error-loop eigenvalues strictly
inside the unit disk.

\emph{Offset-free (observer) side.} The integrating block of
$A_{\text{aug}}$ has eigenvalues at $1$ and is \emph{not} controllable from
$\Fmpc$---the input $[B_d^{(m)};0]$ is zero on the $\dhat$ rows. This is by
design: a constant disturbance is rejected through the internal-model plus
observer, not by driving the $\dhat$ modes. The property actually required is
\emph{detectability} of $(A_{\text{aug}},C)$ with $C=[I_6\;0]$; since $\dhat$
enters the measured state $x_e$ through $B_d^{(m)}$ of full column rank, the
augmented pair is detectable, and the steady-state Kalman filter gives
$\dhat(k)\to d$ for any bounded constant $d$. The free response
\eqref{eq:freeresponse} then pre-loads $-\dhat$, and the stable error loop
converges to $e_\infty=0$. The box constraint $\|\Fmpc\|_\infty\le F_{\max}$
defines a polyhedral invariant set $\mathcal{S}$~\cite{cao2002antiwindup} inside
which this convergence holds.
\end{proof}

\subsection{Transient Bound Across Contact Transitions}

Across a contact switch, $B_d$ jumps by
$\Delta B_d=B_d^{(m_{\text{new}})}-B_d^{(m_{\text{old}})}$.
The error bound is:
\begin{equation}\label{eq:transient}
  \|e(t)\| \leq \|e(t_s)\|
    + c_1\|\Delta B_d\|\cdot\|\dhat\|
    + c_2\|d_{\text{new}}\|\cdot\Delta t_{\text{conv}}
\end{equation}
where $\Delta t_{\text{conv}}\approx5\tau_{\text{Kalman}}$ is the
re-convergence time after covariance inflation, and $c_1,c_2>0$ depend on
the closed-loop eigenvalue placement.

\subsection{Null-Space Barrier Stability}

The contact-consistent null-space torque enforcing joint limits is:
\begin{equation}\label{eq:null}
  \tau_{\text{null}}
    = \bar{N}_{\text{arm}}^\top\bigl(-k_{\text{null}}(q-q_0)
      -d_{\text{null}}\dot{q}+g(q)\bigr)
\end{equation}
where $\bar{N}_{\text{arm}}=I-\bar{J}_{\text{arm}}J_{\text{arm}}$ uses the
contact-consistent pseudoinverse and $g(q)$ is the joint-limit barrier
gradient.
The projection through $\bar{N}_{\text{arm}}$ guarantees zero wrench at the
arm task coordinate.

% ----------------------------------------------------------
\section{Complete Torque Equation and Architectural Comparison}
\label{sec:torque}

\subsection{Complete Joint Torque Equation}

Combining all three levels:
\begin{equation}\label{eq:fulltorque}
  \tau = \tau_{\text{contact}}
       + \bar{N}_1^\top\tau_{\text{balance}}
       + \Naug^\top\bigl[\tau_{\text{ff,arm}}
         + J_{\text{arm}}^\top\Fmpc + \tau_{\text{null}}\bigr]
\end{equation}
The hierarchical null-space structure follows the SK05
law~\cite{sentis2005synthesis} extended to the contact-consistent
floating-base setting~\cite{khatib1987unified,righetti2011inverse}; the MPC
layer $\Fmpc$ and its Kalman augmentation are the contributions of this
paper. The hierarchical null-space projections $\bar{N}_1^\top$ and
$\Naug^\top$ mathematically guarantee that the arm Impedance MPC layer
cannot disturb the contact-maintenance and balance layers above it, while
the Kalman-augmented QP in $\Fmpc$ provides the disturbance rejection that
the classical OS PD law \eqref{eq:PD} cannot achieve.

\subsection{Architectural Comparison}\label{sec:comparison}

Table~\ref{tab:comparison} summarizes the mathematical and architectural
distinctions between the proposed framework and the most relevant prior
work.

\begin{table*}[!t]
\caption{Architectural and Mathematical Comparison}
\label{tab:comparison}
\renewcommand{\arraystretch}{1.3}
\begin{tabular*}{\textwidth}{@{\extracolsep{\fill}}llllll}
\toprule
\textbf{Aspect} & \textbf{Base \cite{cao2026impedance}} & \textbf{Bellicoso et al.\ \cite{bellicoso2016perception}} & \textbf{Kim et al.\ \cite{kim2019highly}} & \textbf{Proposed} \\
\midrule
Primary objective    & Fixed-base pHRI           & Quadruped locomotion  & Dynamic locomotion    & Floating-base pHRI + manipulation \\
Task-space inertia   & $(JM^{-1}J^\top)^{-1}$   & N/A                   & N/A                   & $(J\Mbar J^\top)^{-1}$ \\
Prediction horizon   & $N$-step QP               & $N=1$                 & $\sim$10--30 steps    & $N$-step QP (${\geq}$1\,kHz) \\
Disturbance handling & Kalman, fixed base        & WBC weight tuning     & Centroidal inertia    & Kalman: pHRI + leg momentum \\
Input matrix         & $B_d$ fixed               & N/A                   & N/A                   & $B_d^{(m)}$, contact-mode indexed \\
Null-space use       & Joint centering           & Posture tracking      & Locomotion            & Predictive impedance QP \\
Steady-state error   & Zero (fixed base)         & Nonzero under load    & Nonzero under load    & Zero (Theorem~\ref{thm:zss}) \\
Contact transitions  & N/A                       & Mode switching        & Gait phases           & Covariance-inflation protocol \\
\bottomrule
\end{tabular*}
\end{table*}

% ----------------------------------------------------------
\section{Simulation Benchmarks}\label{sec:simulation}

\subsection{Simulation Platform}

All experiments were conducted in MuJoCo~3.2~\cite{todorov2012mujoco} at a
2\,kHz integration rate.
Scenarios~A and~B use a biped comprising a 3-DOF right arm, two 4-DOF legs,
and a 6-DOF unactuated floating base (17~DOF total, 11~actuated), with
total mass 46\,kg.
Scenario~C uses the official \textbf{Unitree G1} MJCF model from MuJoCo
Menagerie (29~DOF, 33.3\,kg, factory CAD kinematics), augmented with a
single end-effector site at \texttt{right\_wrist\_yaw\_link}.
The G1 model uses position actuators ($K_p=500$, $\text{dampratio}=1$);
Level~3 applies the position-as-torque approximation
($\Delta q_i=\tau_i/K_p$) to inject impedance forces through the position
channel.

The legs are held in double support by a joint-space PD balance controller;
the centroidal-MPC balance planner of Section~\ref{sec:impedancempc} is not
exercised in these static-/fixed-stance scenarios. The arm interaction layer
(Impedance MPC) runs at 1\,kHz with $N=20$, $\Delta t_3=1$\,ms.
Friction-cone half-angle $\mu=0.6$.
Level~3 cost weights: $Q=\mathrm{diag}(6\times10^4 I_3,\;60\,I_3)$,
$R=0.01\,I_3$.
Kalman noise: $Q_w=\mathrm{diag}(10^{-4}I_6,\;10^{-2}I_3)$,
$R_v=10^{-6}I_3$.
Covariance-inflation coefficient $\alpha=4$.

\subsection{Benchmarked Controllers}

Seven controllers are evaluated across all scenarios (Table~\ref{tab:controllers}).

\begin{table}[!t]
\caption{Benchmarked Controllers}
\label{tab:controllers}
\renewcommand{\arraystretch}{1.2}
\begin{tabular}{lp{6cm}}
\toprule
\textbf{Label} & \textbf{Description} \\
\midrule
D1 & Operational-space PD (task PD, no priority hierarchy), $K_P=800$\,N/m, $K_D=40$\,Ns/m \\
D2 & Operational-space PI: adds $K_I=150$\,N/(m$\cdot$s) with anti-windup \\
D3 & Fixed-base impedance MPC~\cite{cao2026impedance} (uses $M^{-1}$ instead of $\Mbar$) \\
D4 & WBC assembler + null-space PD centering (no prediction) \\
D5 & Proposed: WBC + Impedance MPC, no Kalman \\
D6 & Proposed: WBC + Impedance MPC + Kalman, no inflation ($\alpha=1$) \\
D7 & Proposed full: WBC + Impedance MPC + Kalman + covariance inflation ($\alpha=4$) \\
\bottomrule
\end{tabular}
\end{table}

D1 establishes the analytical baseline: for $K_P=800$\,N/m and
$F_h=8$\,N, \eqref{eq:sse} predicts $e_\infty=F_h/K_P=10.0$\,mm exactly.

\subsection{Scenario A: Fixed Double-Support Step Disturbance}

The robot holds a stationary double-support stance.
At $t=0.5$\,s, a step pHRI force of 8\,N is applied at the end-effector in
the $x$-direction and held for the remaining 4.5\,s.
Results are shown in Table~\ref{tab:scenarioA} and Fig.~\ref{fig:scenarioA}.

\begin{table}[!t]
\caption{Scenario A --- Fixed Stance, 8\,N Step Disturbance}
\label{tab:scenarioA}
\renewcommand{\arraystretch}{1.2}
\begin{tabular}{lcc}
\toprule
\textbf{Controller} & \textbf{RMS err [mm]} & \textbf{SS err [mm]} \\
\midrule
D1 OS PD          & 9.24  & 10.17 \\
D2 OS PI          & 6.82  &  5.86 \\
D3 Fixed-base MPC   & 0.21  &  0.15 \\
D4 WBC + PD         & 10.06 & 10.87 \\
D5 Proposed, no Kalman & 0.68 & 0.74 \\
D6 Proposed, no inflation & 0.43 & 0.40 \\
D7 Proposed full    & \textbf{0.43} & \textbf{0.40} \\
\bottomrule
\end{tabular}
\end{table}

\begin{figure}[!t]
  \centering
  \includegraphics[width=\columnwidth]{scenario_a_results.png}
  \caption{Scenario A --- Fixed double-support stance, 8\,N step pHRI at
           $t=0.5$\,s.
           \textit{Top:} Cartesian end-effector error norm $\|e\|$ over
           time.
           In static double support all predictive controllers---including
           the fixed-base D3---reach sub-millimetre steady state; contact
           consistency does not distinguish them here.
           \textit{Bottom:} RMS and steady-state~(SS) error bar chart.
           D4--D6 are omitted for readability and reported in
           Table~\ref{tab:scenarioA}.}
  \label{fig:scenarioA}
\end{figure}

Key findings: (1)~The reactive baselines (D1, D4) leave the theoretical
$10$\,mm impedance offset ($8\,\text{N}/800\,\text{N/m}$). (2)~In static
double support \emph{all} predictive controllers reach sub-millimetre steady
state---including the fixed-base D3 ($0.15$\,mm): with sufficient feedback
gain the offset is suppressed regardless of the inertia model, because a
static stance never excites the contact-mode dynamics. Contact-consistency is
thus \emph{not} what drives static precision (its value appears only under
dynamic loading, Scenario~B). (3)~The Kalman augmentation gives a modest
steady-state improvement (D5 $0.74\to$ D7 $0.40$\,mm, $1.8\times$), and
covariance inflation has no effect here (D6$=$D7) because the contact mode
never switches in this scenario.

\subsection{Scenario B: Stance with Periodic Contact-Transition Shocks}

The robot holds double-support stance while a sustained 8\,N pHRI force
and periodic 6\,N spikes (1\,Hz, 0.1\,s duration) are applied.
Results are shown in Table~\ref{tab:scenarioB} and Fig.~\ref{fig:scenarioB}.

\begin{table}[!t]
\caption{Scenario B --- Stance + 1\,Hz Shocks, Sustained 8\,N pHRI}
\label{tab:scenarioB}
\renewcommand{\arraystretch}{1.2}
\begin{tabular}{lcc}
\toprule
\textbf{Controller} & \textbf{RMS err [mm]} & \textbf{Peak at switch [mm]} \\
\midrule
D1 OS PD          & 10.94 & 16.42 \\
D2 OS PI          &  5.91 & 13.71 \\
D3 Fixed-base MPC   & \textbf{372.3} & \textbf{605.9} \\
D4 WBC + PD         & 11.89 & 17.68 \\
D5 Proposed, no Kalman &  0.82 &  1.32 \\
D6 Proposed, no inflation &  0.50 &  1.15 \\
D7 Proposed full    & \textbf{0.50} & \textbf{1.15} \\
\bottomrule
\end{tabular}
\end{table}

This scenario is the decisive test. Under the periodic shocks the fixed-base
D3---among the best controllers in the static Scenario~A---\emph{diverges}
(RMS $372$\,mm): its free-space inertia under-estimates the true
contact-coupled apparent inertia at the hand (${\sim}5$ vs.\ ${\sim}1.2$\,kg),
producing an over-aggressive predictive gain that the dynamic excitation
destabilizes. The contact-consistent controllers (D5--D7) remain stable and
sub-millimetre (RMS $\leq0.82$\,mm) because $\Larm$ correctly reflects the
higher apparent inertia and hence a well-conditioned gain. \emph{This is where
contact-consistency earns its role}---not in static precision, but in
stability under dynamic contact loading. The Kalman again yields a modest gain
(D5 $0.82\to$ D7 $0.50$\,mm). Covariance inflation is inert here (D6$=$D7):
despite the label, this scenario applies force spikes \emph{without} an actual
foot lift-off or touchdown, so the contact mode never changes and there is
nothing for inflation to re-estimate. A genuine test of the inflation protocol
requires a real contact-set change, provided by the bracing-hand transition of
Scenario~E.

\begin{figure}[!t]
  \centering
  \includegraphics[width=\columnwidth]{scenario_b_results.png}
  \caption{Scenario B --- double-support stance, sustained 8\,N pHRI
           + 6\,N periodic shocks at 1\,Hz.
           The contact-consistent controllers (D5--D7) stay sub-millimetre,
           while the fixed-base D3 diverges under the dynamic excitation
           (Table~\ref{tab:scenarioB}).}
  \label{fig:scenarioB}
\end{figure}

\subsection{Scenario C: Unitree G1 Real Model, Fixed Stance}

Scenario~A is repeated on the official Unitree G1 MJCF (29~DOF, 33.3\,kg).
Results are shown in Table~\ref{tab:scenarioC} and Fig.~\ref{fig:scenarioC}.

\begin{table}[!t]
\caption{Scenario C --- Unitree G1 (33.3\,kg), Fixed Stance, 8\,N Step}
\label{tab:scenarioC}
\renewcommand{\arraystretch}{1.2}
\begin{tabular}{lcc}
\toprule
\textbf{Controller} & \textbf{RMS err [mm]} & \textbf{SS err [mm]} \\
\midrule
D1 OS PD          & 9.00  & 9.70  \\
D2 OS PI          & 6.41  & 5.39  \\
D3 Fixed-base MPC   & 1.38  & 2.28  \\
D4 WBC + PD         & 9.00  & 9.70  \\
D5 Proposed, no Kalman & 1.61 & 2.45 \\
D6 Proposed, no inflation & 1.50 & 2.41 \\
D7 Proposed full    & \textbf{1.50} & \textbf{2.41} \\
\bottomrule
\end{tabular}
\end{table}

\begin{figure}[!t]
  \centering
  \includegraphics[width=\columnwidth]{scenario_c_g1_results.png}
  \caption{Scenario C --- Unitree G1 official model (33.3\,kg, 29~DOF),
           fixed stance, 8\,N step pHRI.
           In static G1 stance the contact-consistent (D7) and fixed-base
           (D3) predictive controllers are comparable; all improve
           ${\sim}4\times$ over the reactive D1.
           D4--D6 are omitted for readability and reported in
           Table~\ref{tab:scenarioC}.}
  \label{fig:scenarioC}
\end{figure}

As in Scenario~A, the static G1 stance does not distinguish contact-consistent
from fixed-base predictive control: D3 ($2.28$\,mm) and D7 ($2.41$\,mm) are
comparable, and all predictive controllers improve ${\sim}4\times$ over the
reactive D1 ($9.70$\,mm). The Kalman benefit is negligible in this static case
(D5 $2.45\to$ D7 $2.41$\,mm). The residual few-millimetre error (vs.\
sub-millimetre on the biped) arises from the 5\,Hz position-actuator bandwidth
attenuating 1\,kHz MPC corrections; direct joint-torque mode on G1/R1~EDU
hardware is expected to recover near-full performance.

\subsection{Scenario E: Bracing-Hand Support Transition}

To exercise a \emph{genuine} contact-mode transition without single-support
balance, the biped is given a left arm that periodically braces against and
releases a fixed rail while both feet remain planted.
The active contact set alternates
$\{\text{L foot},\text{R foot}\}\leftrightarrow\{\text{L foot},\text{R foot},\text{L hand}\}$
($J_c$: $6\leftrightarrow9$ rows), so $\Mbar$, the right-arm task inertia
$\Larm$, and the input matrix $B_d$ genuinely switch at each brace and release.
A sustained 8\,N pHRI force acts on the right (task) arm throughout.

\begin{table}[!t]
\caption{Scenario E --- Bracing-Hand Support Transition, Sustained 8\,N pHRI}
\label{tab:scenarioE}
\renewcommand{\arraystretch}{1.2}
\begin{tabular}{lcc}
\toprule
\textbf{Controller} & \textbf{RMS err [mm]} & \textbf{Peak at switch [mm]} \\
\midrule
D5 Proposed, no Kalman              & 2.80 & 4.50 \\
D6 Kalman, no inflation ($\alpha=1$) & 2.72 & 4.67 \\
D7 Kalman + inflation ($\alpha=4$)   & \textbf{2.70} & \textbf{4.42} \\
\bottomrule
\end{tabular}
\end{table}

This scenario validates the covariance-inflation protocol. At each contact-mode
switch the Kalman input matrix $B_d^{(m)}$ jumps discontinuously; without
inflation (D6) the stale covariance briefly biases $\dhat$, so the peak
transition error ($4.67$\,mm) is in fact \emph{worse} than the no-Kalman
controller (D5, $4.50$\,mm). Inflating the covariance at the switch
(D7, $\alpha=4$) widens the filter, accelerates re-estimation of $\dhat$ under
the new $B_d$, and recovers the lowest peak ($4.42$\,mm). The effect is modest
in magnitude but confirms the protocol's mechanism on a genuine contact-set
change; a full single$\leftrightarrow$double-support transition, which requires
the Level-1 balance controller, remains future work.

% ----------------------------------------------------------
\section{Conclusion}\label{sec:conclusion}

This paper developed the \emph{interaction layer} of a floating-base
whole-body controller: a null-space Impedance MPC that rejects pHRI
disturbances at the arm end-effector with zero steady-state error at fixed
contact mode, within the residual null space of the balance tasks.
The central technical contribution is the proof that a contact-consistent
feedback linearization reduces the arm end-effector plant to a linear double
integrator with a constant discrete state matrix within each contact mode,
enabling offline precomputation of the QP cost inverse and ${\geq}1$\,kHz
operation.
The contact-mode-indexed input-matrix library and a Kalman
covariance-inflation protocol carry these properties across contact
transitions. The inflation protocol is validated on the bracing-hand support
transition of Scenario~E, where the active contact set genuinely switches while
both feet stay planted: inflation recovers the lowest peak transition error,
and without it the Kalman is briefly worse than the no-Kalman controller at each
switch. A full single$\leftrightarrow$double-support transition---which requires
the Level-1 balance controller outside the interaction layer studied here---
remains future work.
The Impedance Equivalence Theorem establishes formal equivalence with
classical operational-space impedance control in the infinite-horizon limit,
and shows that all three closed-loop impedance parameters~$M_{d,\text{eff}}$,
$D_{\text{eff}}$, $K_{\text{eff}}$ adapt automatically through
$\Larm(q)$ without online re-optimization.

The proposed architecture occupies a structural niche not addressed by prior
locomotion-centric frameworks~\cite{dicarlo2018dynamic,kim2019highly,bellicoso2016perception,koolen2016design}:
it deliberately halts the WBC stack after balance constraints are satisfied
and injects a predictive, compliance-controlled manipulation layer into the
residual null space.

In simulation, the decisive benefit of the contact-consistent formulation is
\emph{stability}: under dynamic contact shocks a fixed-base impedance MPC of
comparable gain diverges, while the contact-consistent layer stays
sub-millimetre, because $\Larm$ correctly reflects the higher contact-coupled
apparent inertia and hence a well-conditioned predictive gain. In static
stance the contact-consistent and fixed-base controllers are comparable, and
the Kalman augmentation yields a modest offset-free refinement: the large
steady-state-error reductions attainable at a fixed configuration are
gain-dominated rather than a unique property of the disturbance observer.

The simulations here run the interaction layer above a stabilizing balance
controller; a primary direction for future work is to close the loop with the
full surrounding architecture---a centroidal MPC over the single-rigid-body
dynamics and a priority-driven contact-consistent WBC hierarchy---so that the
balance and interaction layers are co-optimized rather than interfaced.
Further work includes hardware validation on a Unitree R1 or G1~EDU platform
(the low-level SDK torque interface is directly compatible with the proposed
architecture), extension to variable contact modes during dynamic walking,
and integration of force-torque sensor feedback for improved Kalman
convergence at contact transitions.

% ----------------------------------------------------------
\bibliographystyle{IEEEtran}
\bibliography{references}
